\section{A mixed detector block recovers the precision saving}

The preceding logarithm argument makes only one cross test order12.
The full implemented list has nine order-one tests, one order-twelve test,
and six order-forty tests. For each of the three scales, write
\[
 t=A\lambda,\quad\delta=t-1,\quad
 e=(A\delta)^5/F(t),\quad R=(e\delta)^{-1},\quad
 q_m=e^{2m-1}/\delta^3.
\]
The existing arbitrary-order detector lemma applies because, for each
test $D_j$ and positive integer $m$,
\[
 1+R\left(\frac{D_j/\delta^{m+1}}{R^m}\right)^2
       =1+q_mD_j^2.
\]
The inverse-power circuit already supplies $q_{12}=q_1e^{22}$ and
$q_{40}=q_1e^{78}$. No extra inverse-power chain is needed.

Let $A_1$ denote one order-one squared test and $B_{12}$ the sole
order-twelve squared test. Precompute $A_1B_{12}$ once over $K$, and
evaluate in each detector
\[
 (1+q_1A_1)(1+q_{12}B_{12})
  =1+q_1A_1+q_{12}B_{12}+(q_1q_{12})(A_1B_{12}).
\]
Group the remaining eight order-one factors as $4+4$, and the remaining
six order-forty factors as $3+3$. This retains five evaluated $K$-blocks.
The shared coefficient cost remains 34 rational coordinates: two for the
mixed coefficient, twenty for the order-one groups, and twelve for the
order-forty groups. The extra $q_1$ power and removed $q_{40}$ power cancel.
Only the three scalar products $q_1q_{12}$, one per detector, are new.
The $G_9\to G_2$ change saves six coordinates, so its net saving is three.
Together with the nine primary-root savings, this detector lane saves twelve.

\subsection{Actual domain and completeness consumers}

The retained positive-$\lambda$ condition gives $q_1>0$; the even powers
in $q_{12},q_{40}$ make them positive as well. Every factor is therefore
at least one at both real embeddings. All nonzero trace, norm, and
positive pole-input domains used by the earlier guard removals persist.

At a selected contact prime, $v(R)=1$ and $\delta$ is a unit. At the
successful place, $v(D_j)\ge m_j$, so
\[
 v(q_{m_j}D_j^2)=1-2m_j+2v(D_j)\ge1.
\]
Each inverse factor has residue one there. At either place, the
nonconstant term has odd valuation when nonzero. Its factor thus has
valuation zero or negative, and its inverse has residue one or zero.
The reciprocal trace is a unit, with residue one or two. This is the
same local pole-input consumer as before.

The new primary square factor is $1+q_mf(U_i)$. It depends only on the
underlying point, not on the sign of its inverse-root chart. At infinity,
$|\sigma(g)|^2=\sigma(f)$ supplies fixed bounds from the old point bounds.
For each fixed $m$, the weights satisfy $q_m=O(\lambda^{-3})$, and the
product approaches one as before. At the fixed exceptional primes, the
identity $v(g)=v(f)/2$ changes only finite lower bounds. The same prescribed
divisibility makes all terms sufficiently small. The inherited local and
real pole-completeness statements therefore apply after altering those
finite bounds; no stronger prime-pattern hypothesis is required.

The detailed source proof is preserved as
\texttt{detector\_reduction\_proof.txt}. Its independent source checker
reconstructs all 16 aliases, full conjugation of each alternate ratio,
the five primary root squares, the actual shared coefficients and mixed
evaluations, all inverse weights, positivity, and norm/pole-input consumers.
It checks 304 source assertions and two expanded component identities.
